\documentclass[11pt,a4paper]{beamer}
\usetheme{Madrid}
\usepackage[utf8]{inputenc}
\usepackage{amsmath}
\usepackage{amsfonts}
\usepackage{amssymb}
\usepackage{graphicx}
\usepackage{xcolor}
\usepackage{tikz}
\usepackage{booktabs}
\usepackage{array}
\usepackage{colortbl}
\usepackage{multicol}

% ============================================================
% TINY FONT FOR MAXIMUM CONTENT
% ============================================================
\usefonttheme{professionalfonts}

\setbeamerfont{normal text}{size=\scriptsize}
\setbeamerfont{itemize/enumerate body}{size=\scriptsize}
\setbeamerfont{itemize/enumerate subbody}{size=\tiny}
\setbeamerfont{block body}{size=\scriptsize}
\setbeamerfont{block title}{size=\scriptsize}
\setbeamerfont{frametitle}{size=\small}
\setbeamerfont{title}{size=\normalsize}
\setbeamerfont{subtitle}{size=\small}
\setbeamerfont{author}{size=\tiny}
\setbeamerfont{date}{size=\tiny}

\renewcommand{\scriptsize}{\fontsize{6pt}{7pt}\selectfont}
\renewcommand{\footnotesize}{\fontsize{6.5pt}{8pt}\selectfont}
\renewcommand{\small}{\fontsize{7.5pt}{9pt}\selectfont}
\renewcommand{\normalsize}{\fontsize{8.5pt}{10.5pt}\selectfont}

\newcommand{\redflag}[1]{\textcolor{red}{\textbf{[!] #1}}}
\newcommand{\bluenote}[1]{\textcolor{blue}{\textbf{[i] #1}}}
\newcommand{\greennote}[1]{\textcolor{green!50!black}{\textbf{[CORRECT] #1}}}
\newcommand{\examfocus}[1]{\textcolor{orange!80!black}{\textbf{[FOCUS] #1}}}
\newcommand{\trap}[1]{\textcolor{purple}{\textbf{[TRAP] #1}}}
\newcommand{\correct}[1]{\textcolor{green!50!black}{\textbf{[right] #1}}}
\newcommand{\scenario}[1]{\textcolor{blue!70!black}{\textbf{[SCENARIO] #1}}}
\newcommand{\keyterm}[1]{\textcolor{red!80!black}{\textbf{[KEY] #1}}}
\newcommand{\option}[1]{\textcolor{black}{\textbf{#1}}}
\newcommand{\wrong}[1]{\textcolor{gray}{\textbf{[✗] #1}}}

\title[CAMS Exam: Bitcoin \& Crypto Questions]{CAMS Exam: Complete Bitcoin \& Cryptocurrency\\Hard Question Bank}
\author{All 4 Domains | Multi-Option Questions}
\date{}

\setbeamercovered{transparent}
\setbeamertemplate{navigation symbols}{}
\setbeamertemplate{frametitle continuation}{}
\setlength{\parskip}{2pt}
\setlength{\itemsep}{1pt}

\begin{document}

% ============================================================
% TITLE SLIDE
% ============================================================
\begin{frame}
\titlepage
\centering
\tiny 45+ Hard Questions | VASPs | Mixers | Stablecoins | NFTs | DeFi | Mining | MiCA | Travel Rule | Blockchain Tracing
\end{frame}

% ============================================================
% DOMAIN A - QUESTIONS 1-12
% ============================================================

\begin{frame}{DOMAIN A - QUESTION 1: Cryptoasset Laundering Techniques}
\scenario{A criminal network generates \$3.2 million from darknet drug trafficking in Bitcoin. To obscure the origin of funds, they execute the following steps: (1) Split \$3.2 million BTC into 320 separate transactions of 0.1 BTC each; (2) Route each transaction through three different mixing services (including CoinJoin/Wasabi Wallet); (3) Further split output into 6,400 micro-transactions of 0.005 BTC each; (4) Consolidate funds into 32 new wallets over 45 days; (5) Transfer to a VASP in a jurisdiction with no AML crypto regulation; (6) Convert to USDT stablecoin; (7) Withdraw as fiat currency through 50 separate ATM withdrawals under \$1,000 each.}
\vspace{0.2cm}
\textbf{Question: Which money laundering techniques are being used in this scheme? (Select the BEST answer)}
\vspace{0.1cm}
\begin{itemize}
    \item \option{A)} Structuring and trade-based money laundering only
    \item \option{B)} Microstructuring, mixing/tumbling, and stablecoin conversion
    \item \option{C)} Smurfing and real estate investment only
    \item \option{D)} Cash smuggling and shell company formation only
    \item \option{E)} Microstructuring, mixing, stablecoin conversion, and structured withdrawals
\end{itemize}
\vspace{0.2cm}
\trap{Only microstructuring is happening; mixing is not money laundering.}
\correct{Answer: E. The scheme uses microstructuring (tiny crypto transactions), mixing/tumbling (CoinJoin), stablecoin conversion (USDT to avoid volatility), and structured cash withdrawals (under \$1,000 to avoid CTR triggers). Multiple techniques layered together.}
\redflag{Microstructuring is specifically used with digital asset laundering. Mixers make transaction tracing nearly impossible. Stablecoins at unregulated VASPs are high risk.}
\end{frame}

\begin{frame}{DOMAIN A - QUESTION 2: Privacy Coins Risk Classification}
\scenario{A bank has a customer who receives monthly payments in Monero (XMR), a privacy coin that uses RingCT and stealth addresses to obscure transaction details. The customer also holds a separate account for business operations in Bitcoin and Ethereum. The customer refuses to provide source of funds documentation for the Monero payments, stating "privacy is my right." The bank's transaction monitoring system can trace the Bitcoin and Ethereum transactions but cannot trace the origin or destination of the Monero transactions. The customer's declared annual income is \$80,000, but Monero deposits total approximately \$25,000 per month. The customer has been a client for 3 years with no previous issues.}
\vspace{0.2cm}
\textbf{Question: How should the bank classify the money laundering risk associated with this customer's use of Monero?}
\vspace{0.1cm}
\begin{itemize}
    \item \option{A)} Lower risk than Bitcoin because Monero has stronger security features
    \item \option{B)} Same risk as Bitcoin and Ethereum because all cryptocurrencies have identical AML risk profiles
    \item \option{C)} Heightened risk due to anonymity features and inability to trace transactions, requiring EDD
    \item \option{D)} No risk assessment needed because privacy coins are legal in most jurisdictions
    \item \option{E)} Lower risk because privacy coins are less liquid and harder to convert to fiat
\end{itemize}
\vspace{0.2cm}
\trap{All cryptoassets have the same risk level; privacy is irrelevant.}
\correct{Answer: C. Privacy coins with nonpublic blockchains represent HIGHER risk because they complicate efforts to attribute transactions. The inability to trace source of funds, combined with income inconsistency (\$80k declared vs \$300k annual deposits), requires EDD and potential SAR filing.}
\redflag{Privacy coins (Monero, Zcash, Dash with PrivateSend) are HIGHER risk than transparent blockchains like Bitcoin.}
\end{frame}

\begin{frame}{DOMAIN A - QUESTION 3: Stablecoin Money Laundering Vectors}
\scenario{A financial institution identifies a pattern of activity involving a corporate customer that operates a small retail business with annual revenue of approximately \$500,000. Over a 6-month period, the customer has:
- Purchased \$1.2 million in USDT (Tether) from a VASP with minimal KYC
- Transferred the USDT to a decentralized exchange (DEX) with no KYC requirements
- Swapped USDT for DAI stablecoin through a series of 200 small transactions
- Transferred DAI to a VASP in a jurisdiction with weak AML enforcement
- Cashed out \$1.15 million to a bank account in a country with strict bank secrecy laws
The customer cannot explain the source of the original \$1.2 million, which significantly exceeds business revenue.}
\vspace{0.2cm}
\textbf{Question: What is the primary money laundering vulnerability with fiat-collateralized stablecoins like USDT and USDC in this scenario?}
\vspace{0.1cm}
\begin{itemize}
    \item \option{A)} Stablecoins are impossible to trace on any blockchain
    \item \option{B)} Ease of conversion back into traditional fiat combined with regulatory gaps at unregulated VASPs
    \item \option{C)} Stablecoins cannot be used for cross-border transactions
    \item \option{D)} Stablecoins are always fully backed 1:1 with audited reserves, making them low risk
    \item \option{E)} Only algorithmic stablecoins present money laundering risks
\end{itemize}
\vspace{0.2cm}
\trap{Stablecoins are fully backed and transparent, so low risk.}
\correct{Answer: B. The primary vulnerability is "ease of conversion back into traditional fiat and regulatory gaps," making them attractive for money laundering. The customer exploited unregulated VASPs and jurisdictions with weak enforcement to convert crypto to clean fiat.}
\redflag{Stablecoins at unregulated VASPs = high risk. Always verify VASP licensing and jurisdiction.}
\end{frame}

\begin{frame}{DOMAIN A - QUESTION 4: NFT Wash Trading for Money Laundering}
\scenario{A criminal organization creates 50 unique NFT artworks on a blockchain platform with no KYC requirements. The organization then:
- Lists each NFT for prices ranging from 50 ETH to 500 ETH (\$150,000 to \$1.5 million)
- Uses 20 different shell company wallets to "purchase" the NFTs from themselves
- Funds the purchases using Bitcoin traced to darknet market sales
- Receives ETH from the "sales" into fresh wallets not linked to original illicit funds
- Converts ETH to USDC through a decentralized exchange
- Cashes out USDC through a licensed VASP that only performs basic KYC
The total value laundered through this scheme is approximately \$8.5 million.}
\vspace{0.2cm}
\textbf{Question: What money laundering technique does this NFT scheme represent, and what are the red flags?}
\vspace{0.1cm}
\begin{itemize}
    \item \option{A)} Trade-based money laundering; red flags include inconsistent pricing
    \item \option{B)} NFT wash trading/self-laundering; red flags include overpricing, same entity buying and selling, and anonymous platforms
    \item \option{C)} Traditional structuring; red flags include multiple small transactions
    \item \option{D)} Cash smuggling; red flags include physical transportation of currency
    \item \option{E)} Shell company formation only; no money laundering involved
\end{itemize}
\vspace{0.2cm}
\trap{This is legitimate art trading; NFTs have value determined by market.}
\correct{Answer: B. This is NFT wash trading/self-laundering. Red flags include: (1) overpricing compared to market value, (2) same beneficial owner buying and selling, (3) anonymous platforms with no KYC, (4) use of shell company wallets, (5) funds traced to darknet markets.}
\redflag{NFT wash trading = same person buys own NFT to launder money. Overpricing is the primary red flag.}
\end{frame}

\begin{frame}{DOMAIN A - QUESTION 5: Crypto Mixer/Tumbler Red Flags}
\scenario{A VASP receives the following transaction patterns over a 30-day period:
- Transaction 1: 2.5 BTC from address ending in ...abc, processed through ChipMixer
- Transaction 2: 1.8 BTC from address ending in ...def, processed through Blender.io
- Transaction 3: 3.2 BTC from address ending in ...ghi, showing 15 hops through Wasabi Wallet's CoinJoin
- Transaction 4: 0.5 BTC from address ending in ...jkl, no mixing detected
- Transaction 5: 4.0 BTC from address ending in ...mno, processed through Tornado Cash
The customer's declared business is a small coffee shop with annual revenue of \$150,000. The customer cannot explain the source of the BTC or why mixing services were used.}
\vspace{0.2cm}
\textbf{Question: Which transactions should trigger immediate EDD and potential SAR filing, and why?}
\vspace{0.1cm}
\begin{itemize}
    \item \option{A)} Transaction 4 only, because small transactions are always suspicious
    \item \option{B)} All transactions except Transaction 4; mixing services obscure fund origin and make tracing nearly impossible
    \item \option{C)} None of the transactions; mixing services are legal and privacy is a right
    \item \option{D)} Only Transaction 3 because CoinJoin is the only illegal mixing service
    \item \option{E)} Only Transaction 5 because Tornado Cash is sanctioned by OFAC
\end{itemize}
\vspace{0.2cm}
\trap{Only sanctioned mixers like Tornado Cash are problematic.}
\correct{Answer: B. When mixing services are used, fund tracing becomes "extremely difficult and sometimes impossible." Any funds from known mixers require EDD. The customer cannot explain source or reason for mixing. All mixed transactions (1,2,3,5) require SAR filing.}
\redflag{Mixers include ChipMixer, Blender.io, Wasabi Wallet/CoinJoin, Tornado Cash. Any detected mixer usage = high risk.}
\end{frame}

\begin{frame}{DOMAIN A - QUESTION 6: CBDC AML Advantages vs Cash}
\scenario{A central bank is considering launching a retail CBDC. Two policymakers disagree:
- Policymaker A argues that cash should remain dominant because it provides financial privacy and anonymity
- Policymaker B argues that a CBDC would strengthen AML/CFT efforts
Current cash usage in the country is \$200 billion annually. Estimated tax evasion through undeclared cash transactions is \$15 billion per year. Estimated money laundering through cash-intensive businesses is \$8 billion per year. Terrorism financing investigations are frequently hindered by inability to trace cash donations to unknown sources.}
\vspace{0.2cm}
\textbf{Question: What AML/CFT advantages does a CBDC have over physical cash? (Select the BEST answer)}
\vspace{0.1cm}
\begin{itemize}
    \item \option{A)} CBDC cannot be used for any illegal activity whatsoever
    \item \option{B)} Real-time monitoring of transactions and greater transparency for authorities
    \item \option{C)} CBDC has no AML/CFT advantages over physical cash
    \item \option{D)} CBDC eliminates all financial crime completely
    \item \option{E)} CBDC only helps with terrorism financing, not money laundering
\end{itemize}
\vspace{0.2cm}
\trap{Cash and CBDC are the same; no AML benefits.}
\correct{Answer: B. "Unlike cash, CBDCs can be monitored in real time, providing authorities with greater transparency over transactions. This monitoring can help combat money laundering, tax evasion, and terrorist financing."}
\redflag{CBDC transparency is a major AML advantage over anonymous physical cash.}
\end{frame}

\begin{frame}{DOMAIN A - QUESTION 7: Terrorist Financing via Crypto}
\scenario{A terrorist organization uses the following financial network:
- Legitimate donations collected through a charity front organization in Country A
- Donations in Bitcoin and Monero from sympathizers in 15 countries
- Funds transferred through 4 different mixing services over 8 months
- Cryptocurrency converted to USDT through unregulated VASPs
- USDT sent to wallets controlled by terrorist operatives in conflict zone
- Operatives convert USDT to local currency through informal hawala brokers
- Total funds moved: approximately \$4.5 million over 2 years
- No single transaction exceeds \$3,000
The charity files annual reports showing only \$200,000 in donations, all from traceable bank transfers.}
\vspace{0.2cm}
\textbf{Question: What red flags indicate terrorist financing rather than traditional money laundering?}
\vspace{0.1cm}
\begin{itemize}
    \item \option{A)} Use of cryptocurrency for any transaction
    \item \option{B)} High transaction volumes only
    \item \option{C)} Use of legitimate-looking charities, small transaction sizes to avoid detection, mix of legitimate and illicit funds, and informal value transfer systems (hawala)
    \item \option{D)} Only large cash transactions are red flags
    \item \option{E)} Use of stablecoins over Bitcoin
\end{itemize}
\vspace{0.2cm}
\trap{Terrorist financing only uses illegal funds like drug money.}
\correct{Answer: C. Terrorist financing differs from money laundering: (1) source may be legitimate, (2) funds often move through charities/NGOs, (3) transaction sizes are often small to avoid detection, (4) hawala/ARS frequently used. The charity front reporting \$200k but receiving \$4.5M is a major red flag.}
\redflag{Terrorist financing can use LEGITIMATE funds diverted for illicit purposes, making detection more complex.}
\end{frame}

\begin{frame}{DOMAIN A - QUESTION 8: DeFi Protocol Vulnerabilities}
\scenario{A decentralized finance (DeFi) protocol allows users to:
- Lend and borrow cryptoassets without identity verification
- Swap between 200+ different tokens including privacy coins
- Bridge assets across 12 different blockchains
- Earn yield through liquidity pools with no source of funds checks
- Create new tokens with custom names and supply
- No KYC required; only a wallet address needed
Law enforcement identifies that \$50 million in ransomware proceeds has been laundered through this protocol over 18 months. The protocol developers remain anonymous and cannot be identified. The protocol's governance tokens are widely distributed with no single entity holding >2\%.}
\vspace{0.2cm}
\textbf{Question: What are the primary AML vulnerabilities of this DeFi protocol, and who bears responsibility under FATF guidance?}
\vspace{0.1cm}
\begin{itemize}
    \item \option{A)} No vulnerabilities; DeFi is fully compliant by design
    \item \option{B)} Anonymity, cross-chain bridging, no KYC; the developers may be considered VASPs if they maintain control or influence
    \item \option{C)} Only the bridging feature is a vulnerability
    \item \option{D)} Liquidity pools are the only vulnerability; users bear responsibility
    \item \option{E)} No responsibility applies to DeFi because there is no central entity
\end{itemize}
\vspace{0.2cm}
\trap{DeFi is decentralized so no one is responsible.}
\correct{Answer: B. DeFi vulnerabilities include anonymity, cross-chain transfers, and no KYC. Under FATF guidance, persons who maintain "control or sufficient influence" over a DeFi protocol may be considered VASPs. Factors include governance token ownership, admin keys, ability to upgrade contracts, and fee collection.}
\redflag{DeFi with no KYC + cross-chain bridges + privacy coins = highest ML risk.}
\end{frame}

\begin{frame}{DOMAIN A - QUESTION 9: Crypto Mining AML Obligations}
\scenario{A commercial cryptocurrency mining operation has the following characteristics:
- 5,000 ASIC miners running 24/7
- Monthly electricity cost: \$250,000
- Monthly Bitcoin mined: approximately 15 BTC (current value \$450,000)
- Revenue sources: block rewards (12.5 BTC per block found) + transaction fees
- Mining pool payouts distributed daily to 10 different wallet addresses
- BTC sold monthly through three different VASPs to pay operating expenses
- Some BTC sold peer-to-peer through LocalBitcoins with no KYC
- Mining operation registered as a corporation but not registered as MSB or VASP
- Mining operation claims no AML obligations because "mining is not a financial service"}
\vspace{0.2cm}
\textbf{Question: Does this mining operation have AML/CFT obligations, and why?}
\vspace{0.1cm}
\begin{itemize}
    \item \option{A)} No, mining is purely computational with no financial activity
    \item \option{B)} Yes, when converting BTC to fiat through VASPs and P2P, the miner may need to register as MSB/VASP depending on jurisdiction
    \item \option{C)} Only if mining illegally or without license
    \item \option{D)} Only for transaction fees, not block rewards
    \item \option{E)} No, miners are explicitly exempt from all AML regulations globally
\end{itemize}
\vspace{0.2cm}
\trap{Mining is industrial activity, not financial services.}
\correct{Answer: B. Miners receive newly created cryptoassets and transaction fees. When selling BTC to VASPs or through P2P for fiat, the miner is engaging in money transmission or virtual asset exchange. Many jurisdictions require MSB/VASP registration for commercial mining operations that convert crypto to fiat.}
\redflag{Commercial miners selling crypto for fiat = likely MSB/VASP obligations. P2P sales with no KYC = red flag.}
\end{frame}

\begin{frame}{DOMAIN A - QUESTION 10: Algorithmic Stablecoin Collapse ML Risk}
\scenario{A criminal organization exploits an algorithmic stablecoin (similar to UST/Terra) by:
- Purchasing \$10 million of the stablecoin at \$1.00 using illicit funds
- Executing a coordinated attack on the stablecoin's peg using flash loans
- Stablecoin value collapses to \$0.10 within 48 hours
- Criminal organization sells stablecoin at \$0.10, recovering only \$1 million
- Claims \$9 million loss on tax returns as "capital loss"
- Original \$10 million illicit funds now appear as "legitimate trading losses"
- Organization also short-sold the stablecoin before collapse, earning \$15 million in profits on a separate platform
Net position: \$10 million illicit in, \$16 million clean out (\$1 million recovery + \$15 million short profit)}
\vspace{0.2cm}
\textbf{Question: What money laundering technique is demonstrated by this algorithmic stablecoin exploitation?}
\vspace{0.1cm}
\begin{itemize}
    \item \option{A)} Trade-based money laundering only
    \item \option{B)} Loss harvesting / tax fraud disguised as trading losses, combined with market manipulation profits
    \item \option{C)} Traditional structuring
    \item \option{D)} Cash smuggling
    \item \option{E)} Real estate investment
\end{itemize}
\vspace{0.2cm}
\trap{This is just a bad investment; no money laundering.}
\correct{Answer: B. The organization launders funds by: (1) converting illicit funds to stablecoin, (2) exploiting volatility to generate "legitimate" trading losses for tax purposes, (3) earning clean profits through short positions, (4) claiming tax benefits on the fabricated losses. Algorithmic stablecoins' value fluctuations can "disguise illicit asset transfers and complicate forensic tracking."}
\redflag{Algorithmic stablecoin volatility can be exploited for loss harvesting and tax fraud schemes.}
\end{frame}

\begin{frame}{DOMAIN A - QUESTION 11: Virtual Asset Service Provider Definition}
\scenario{A company offers the following services:
- Cryptocurrency exchange: Bitcoin to Ethereum, no fiat accepted
- Non-custodial wallet where users control their own private keys
- Crypto-to-crypto trading pairs only
- No money transmission services
- No fiat currency on-ramp or off-ramp
- Users identified only by email address; no ID verification
- Company located in a jurisdiction with no crypto regulations
- Annual trading volume: \$500 million
The company claims it is not a financial institution and has no AML obligations because it never touches fiat currency.}
\vspace{0.2cm}
\textbf{Question: Under FATF standards, is this company a Virtual Asset Service Provider (VASP) subject to AML/CFT obligations?}
\vspace{0.1cm}
\begin{itemize}
    \item \option{A)} No, only companies that handle fiat currency are VASPs
    \item \option{B)} Yes, crypto-to-crypto exchanges are VASPs regardless of fiat involvement
    \item \option{C)} Only if the company holds customer funds in custody
    \item \option{D)} Only if trading volume exceeds \$1 billion annually
    \item \option{E)} No, non-custodial wallets are never VASPs
\end{itemize}
\vspace{0.2cm}
\trap{Crypto-to-crypto is outside AML scope because no fiat.}
\correct{Answer: B. VASPs include cryptocurrency exchanges regardless of whether they handle fiat. The definition includes "exchange between virtual assets" (crypto-to-crypto). The lack of customer identification and location in a jurisdiction with no regulations does not exempt the company from FATF standards.}
\redflag{VASP definition includes crypto-to-crypto exchanges, wallet providers, and other virtual asset service providers.}
\end{frame}

\begin{frame}{DOMAIN A - QUESTION 12: Cross-Chain Bridge ML Vulnerability}
\scenario{A criminal organization uses cross-chain bridges to launder funds:
- Initial funds: 500 BTC (\$15 million) from ransomware attacks
- Bridge 1: BTC to renBTC on Ethereum (no KYC at bridge)
- Bridge 2: renBTC to USDC on Solana via Wormhole bridge
- Bridge 3: USDC to BNB on Binance Smart Chain via Multichain
- Bridge 4: BNB to AVAX on Avalanche via LayerZero
- Final: AVAX swapped to XMR (Monero) on a DEX with no KYC
- Total bridges used: 7 different protocols
- Total blockchains traversed: 8 different networks
- Time elapsed: 3 hours
- Tracing cost for law enforcement: approximately \$500,000 and 6 months
The organization successfully cashes out \$13.5 million through 20 different VASPs with minimal KYC.}
\vspace{0.2cm}
\textbf{Question: What AML vulnerability do cross-chain bridges present?}
\vspace{0.1cm}
\begin{itemize}
    \item \option{A)} Bridges are fully transparent and easy to trace
    \item \option{B)} Bridges create complex transaction paths across multiple blockchains, making fund tracing extremely difficult and costly for investigators
    \item \option{C)} Bridges only work for stablecoins, not Bitcoin
    \item \option{D)} All bridges require KYC, eliminating ML risk
    \item \option{E)} Bridges are regulated by central banks in all jurisdictions
\end{itemize}
\vspace{0.2cm}
\trap{Bridges are simple and traceable like normal blockchain transactions.}
\correct{Answer: B. Cross-chain bridges allow funds to move between otherwise separate blockchains. Each bridge crossing creates a break in the forensic trail. Multiple bridges across multiple chains (Bitcoin → Ethereum → Solana → BSC → Avalanche) exponentially increases tracing difficulty.}
\redflag{Cross-chain bridges = major forensic challenge. Each bridge crossing breaks the transaction trail.}
\end{frame}

% ============================================================
% DOMAIN B - QUESTIONS 13-24
% ============================================================

\begin{frame}{DOMAIN B - QUESTION 13: FATF Travel Rule for Crypto}
\scenario{A VASP receives a \$25,000 USDC transfer from another VASP. The incoming transfer message contains:
- Originator name: "CryptoTrader88" (username, not legal name)
- Originator address: blank
- Originator account number: wallet address only
- Beneficiary name: complete and verified
- Beneficiary address: complete and verified
- Transaction hash: provided
The sending VASP is located in a jurisdiction that has not implemented the Travel Rule. The receiving VASP is in a jurisdiction that has fully implemented Travel Rule requirements.}
\vspace{0.2cm}
\textbf{Question: Under the FATF Travel Rule, is this transaction compliant, and what should the receiving VASP do?}
\vspace{0.1cm}
\begin{itemize}
    \item \option{A)} Fully compliant because transaction hash is provided
    \item \option{B)} Not compliant; missing required originator information (legal name and address); receiving VASP should reject or file SAR
    \item \option{C)} Compliant because sending VASP's jurisdiction has not implemented Travel Rule
    \item \option{D)} Only need originator name; address is optional for crypto
    \item \option{E)} Travel Rule does not apply to cryptoasset transfers
\end{itemize}
\vspace{0.2cm}
\trap{Travel Rule only applies to banks, not crypto.}
\correct{Answer: B. The Travel Rule applies to VASPs same as banks. Required information: originator name, originator address, originator account number, beneficiary name, beneficiary address. Username is insufficient. Missing information should trigger rejection or SAR filing.}
\redflag{Travel Rule applies to VASPs. Required: full legal name and address, not just username or wallet address.}
\end{frame}

\begin{frame}{DOMAIN B - QUESTION 14: DeFi Protocol VASP Status Determination}
\scenario{A DeFi lending protocol has the following characteristics:
- Governance token distribution: Founder A (35\%), Founder B (28\%), Founder C (12\%), Public (25\%)
- Multi-signature wallet requires 3 of 5 signatures to upgrade smart contracts
- Signers: Founder A, Founder B, Founder C, and two anonymous community members
- Protocol collects 0.3\% fee on all loans (annual fees: \$15 million)
- Fees distributed to treasury controlled by multi-sig wallet
- Protocol marketed by Founders A, B, and C at conferences and on social media
- No KYC required to use the protocol
- Legal entity: None; protocol claims to be "fully decentralized"
Law enforcement identifies that \$200 million in stolen funds has been laundered through the protocol.}
\vspace{0.2cm}
\textbf{Question: Under FATF guidance on DeFi, which parties are likely considered VASPs with AML obligations?}
\vspace{0.1cm}
\begin{itemize}
    \item \option{A)} No one; DeFi is fully decentralized with no responsible party
    \item \option{B)} Only the anonymous community members are VASPs
    \item \option{C)} Founders A, B, and C (control via governance tokens and multi-sig, collect fees, market the protocol)
    \item \option{D)} Only users of the protocol are VASPs
    \item \option{E)} Only the public token holders are VASPs
\end{itemize}
\vspace{0.2cm}
\trap{DeFi is decentralized; no one has responsibility.}
\correct{Answer: C. FATF guidance: persons with "control or sufficient influence" are VASPs. Indicators here: (1) governance token ownership >50\% combined, (2) multi-sig control, (3) fee collection, (4) marketing the protocol to users. Founders A, B, and C must register and implement AML controls.}
\redflag{Governance tokens + admin keys + fee collection + marketing = VASP status regardless of "decentralized" claims.}
\end{frame}

\begin{frame}{DOMAIN B - QUESTION 15: Blockchain Analytics Investigation}
\scenario{A law enforcement agency is investigating a ransomware attack where the victim paid 200 BTC (approximately \$6 million) to a wallet address. The investigator uses blockchain analytics tools and finds:
- The initial receiving address received funds from 50 different addresses
- Those 50 addresses received from 500 addresses
- Those 500 addresses received from 5,000 addresses
- Two addresses among the 5,000 are linked to a regulated VASP through KYC data
- One address shows activity consistent with an OTC trading desk
- Transaction timestamps show patterns consistent with automated laundering (30-minute intervals, same amounts)
- Some addresses received funds from known mixing services
The investigator needs to identify the real-world person behind the ransomware.}
\vspace{0.2cm}
\textbf{Question: What blockchain analytics techniques should the investigator use to identify the suspect?}
\vspace{0.1cm}
\begin{itemize}
    \item \option{A)} Only look at the initial receiving address; tracing further is impossible
    \item \option{B)} On-chain analytics to trace transaction flows and off-chain analytics to link addresses to real-world identities through VASP KYC data and OTC desk records
    \item \option{C)} Only use off-chain analytics; on-chain is useless
    \item \option{D)} Wait for the suspect to confess
    \item \option{E)} Only analyze timestamps; address data is irrelevant
\end{itemize}
\vspace{0.2cm}
\trap{Bitcoin is anonymous; cannot identify anyone.}
\correct{Answer: B. On-chain analytics examine transaction data, wallet addresses, and patterns. Off-chain analytics integrate external sources: VASP KYC data, regulatory filings, social media, OTC desk records. The addresses linked to regulated VASPs provide the link to real-world identities.}
\redflag{On-chain + off-chain analytics = identify real-world persons behind blockchain addresses.}
\end{frame}

\begin{frame}{DOMAIN B - QUESTION 16: Crypto ATM AML Controls}
\scenario{A cryptocurrency ATM operator has 500 machines located in convenience stores, gas stations, and shopping malls across 3 countries. The ATMs allow users to:
- Buy Bitcoin with cash (no ID required for transactions under \$900)
- Sell Bitcoin for cash (no ID required for transactions under \$900)
- Transaction limits: \$5,000 per day, \$25,000 per month
- Daily transaction volume across all machines: approximately \$2.5 million
- Operator claims machines are "just vending machines, not financial services"
- Operator has not registered as MSB/VASP in any jurisdiction
- Operator conducts no KYC, no transaction monitoring, no SAR filing
- Law enforcement finds that 40\% of transactions are linked to known scam wallets}
\vspace{0.2cm}
\textbf{Question: Under FATF standards, does this crypto ATM operator have AML/CFT obligations?}
\vspace{0.1cm}
\begin{itemize}
    \item \option{A)} No, ATMs are automated machines with no human involvement
    \item \option{B)} Yes, the operator is a VASP/MSB with full AML obligations including registration, KYC, transaction monitoring, and SAR filing
    \item \option{C)} Only for transactions over \$10,000
    \item \option{D)} Only if the operator also runs an exchange
    \item \option{E)} No, cash transactions are exempt from crypto regulations
\end{itemize}
\vspace{0.2cm}
\trap{Crypto ATMs are unregulated; no AML needed.}
\correct{Answer: B. Crypto ATMs are VASPs/MSBs. Obligations include: registration/licensing, KYC/CDD (including thresholds lower than \$900 in many jurisdictions), transaction monitoring, suspicious activity reporting, sanctions screening. The \$900 threshold without ID is below many regulatory requirements.}
\redflag{Crypto ATMs = VASPs. Low or no ID thresholds = major red flag for regulators.}
\end{frame}

\begin{frame}{DOMAIN B - QUESTION 17: Lightning Network AML Challenges}
\scenario{A criminal organization begins using the Bitcoin Lightning Network to launder funds. The Lightning Network has the following characteristics:
- Payment channels allow off-chain transactions between parties
- Only channel opening and closing transactions appear on the Bitcoin blockchain
- Millions of payments can occur within a channel without on-chain records
- Payments can be routed through multiple nodes with no public ledger
- Channel balances are known only to channel participants
- Typical transaction size: small to medium ($5 to $500)
- Settlement time: milliseconds to seconds
The organization opens a single channel with \$500,000 and then routes 50,000 small payments through 1,000 different nodes over 6 months. Only the initial \$500,000 channel opening and final channel closing appear on-chain.}
\vspace{0.2cm}
\textbf{Question: What AML challenges does the Lightning Network present for investigators?}
\vspace{0.1cm}
\begin{itemize}
    \item \option{A)} No challenges; Lightning is fully transparent
    \item \option{B)} Only on-chain transactions are traceable; millions of off-chain Lightning payments leave no public record and cannot be traced
    \item \option{C)} Only large payments are hidden; small payments are traceable
    \item \option{D)} All Lightning payments appear on-chain within 24 hours
    \item \option{E)} Lightning Network requires KYC for all users
\end{itemize}
\vspace{0.2cm}
\trap{Lightning is just faster Bitcoin; same traceability.}
\correct{Answer: B. The Lightning Network creates a significant AML challenge because millions of off-chain payments are not recorded on the public blockchain. Only channel opening and closing transactions are visible to investigators. The 50,000 internal payments cannot be traced, and routing paths are not publicly disclosed.}
\redflag{Lightning Network = off-chain payments with no public record. Major forensic blind spot.}
\end{frame}

\begin{frame}{DOMAIN B - QUESTION 18: Non-Custodial Wallet Risk Assessment}
\scenario{A bank has a customer who uses a non-custodial (self-custody) wallet such as MetaMask or Ledger. The customer:
- Receives regular crypto deposits from 50 different wallet addresses
- Sends crypto to 200 different addresses per month
- Also maintains a traditional bank account with the institution
- Transfers fiat to the bank account from a VASP
- Cannot explain the source of crypto funds
- Claims "my wallet is private, I don't have to tell you"
- Transaction volume: approximately \$100,000 per month in crypto
- Bank can see fiat transfers from VASP but cannot see the underlying crypto transactions
The bank's AML system only monitors fiat transactions, not crypto wallet activity.}
\vspace{0.2cm}
\textbf{Question: What is the bank's AML obligation regarding this customer's non-custodial wallet activity?}
\vspace{0.1cm}
\begin{itemize}
    \item \option{A)} No obligation; non-custodial wallets are outside bank's scope
    \item \option{B)} Bank must conduct EDD on source of funds, request wallet transaction history, and file SAR if customer refuses to provide information
    \item \option{C)} Bank can ignore because crypto is unregulated
    \item \option{D)} Only monitor fiat; crypto activity is irrelevant
    \item \option{E)} Close account immediately without investigation
\end{itemize}
\vspace{0.2cm}
\trap{Bank only responsible for fiat transactions, not crypto.}
\correct{Answer: B. The bank has a fiat relationship with the customer. The customer's crypto activity is relevant to understanding source of funds and risk. The bank must conduct EDD. Customer's refusal to explain source of funds is a red flag requiring SAR filing. The bank cannot ignore crypto activity just because it occurs in non-custodial wallets.}
\redflag{Non-custodial wallets = no built-in KYC. Bank must still assess risk of fiat on/off ramps.}
\end{frame}

\begin{frame}{DOMAIN B - QUESTION 19: Crypto Derivatives ML Risk}
\scenario{A customer trades crypto derivatives (futures, options, perpetual swaps) on an unregulated offshore exchange. The customer:
- Deposits 100 BTC (\$3 million) into exchange
- Executes 5,000 trades over 3 months
- Losses: \$2.5 million; Gains: \$1.5 million; Net loss: \$1 million
- Withdraws 50 BTC (\$1.5 million) after 3 months
- Claims the \$1.5 million withdrawal is "trading profits" despite net loss
- Original 100 BTC source: unknown
- Exchange has no KYC; customer used VPN to access
- Customer's declared income: \$120,000/year from restaurant business}
\vspace{0.2cm}
\textbf{Question: What money laundering technique might be occurring through crypto derivatives trading?}
\vspace{0.1cm}
\begin{itemize}
    \item \option{A)} No money laundering; trading losses prove legitimate activity
    \item \option{B)} Loss harvesting to legitimize withdrawals, layering through complex trades, and commingling illicit funds with "trading activity" to obscure source
    \item \option{C)} Only derivatives with fiat settlement are problematic
    \item \option{D)} All crypto derivatives are fully traceable and low risk
    \item \option{E)} Only options trading, not futures
\end{itemize}
\vspace{0.2cm}
\trap{Trading losses mean no profit; nothing to launder.}
\correct{Answer: B. Even with net losses, criminals can: (1) commingle illicit funds with trading activity, (2) use complex trades to layer funds, (3) claim withdrawals as "legitimate trading proceeds," (4) exploit complex derivative products to obscure fund origin. The withdrawal amount being claimed as "profits" despite net loss is suspicious.}
\redflag{Crypto derivatives = complex layering opportunity. Losses don't mean no laundering occurred.}
\end{frame}

\begin{frame}{DOMAIN B - QUESTION 20: VASP Licensing Jurisdiction Analysis}
\scenario{A VASP operates in the following manner:
- Legal entity incorporated in Country A (no crypto regulations)
- Servers located in Country B (cloud provider)
- Management team located in Country C, D, and E
- Customers from 50 countries
- Bank accounts in Country F and G
- Payment processing through Country H
- No physical office anywhere
- Claims "we are registered nowhere because we are everywhere"
- Annual trading volume: \$5 billion
- Law enforcement from 15 countries has requested customer data; VASP refuses all requests}
\vspace{0.2cm}
\textbf{Question: Under FATF standards, which jurisdictions have responsibility for regulating this VASP?}
\vspace{0.1cm}
\begin{itemize}
    \item \option{A)} No jurisdiction has responsibility
    \item \option{B)} Country A (incorporation), Country F/G (banking), and any jurisdiction where the VASP has meaningful presence or customers
    \item \option{C)} Only Country B where servers are located
    \item \option{D)} Only the customer's home countries
    \item \option{E)} No regulation needed for crypto
\end{itemize}
\vspace{0.2cm}
\trap{No physical presence = no regulation anywhere.}
\correct{Answer: B. FATF expects VASPs to be licensed/registered in the jurisdiction where they are incorporated, where they have their place of management, and where they have significant customer base. The VASP's attempt to be "nowhere" is a red flag. Multiple jurisdictions may claim responsibility. Such VASPs are often designated as high-risk.}
\redflag{VASPs with no clear jurisdiction = high risk. FATF expects licensing where incorporated and where management located.}
\end{frame}

\begin{frame}{DOMAIN B - QUESTION 21: CBDC Offline Payment AML Controls}
\scenario{A central bank is designing a retail CBDC with both online and offline capabilities. The offline functionality would allow payments without internet connectivity using hardware wallets (similar to prepaid cards). The central bank proposes:
- Offline transaction limit: \$200 per transaction
- Offline daily limit: \$500
- Offline wallet maximum balance: \$1,000
- No real-time monitoring for offline transactions
- Offline wallets can be "topped up" online
- No KYC required for offline wallets under \$500 balance
- Annual offline transaction volume estimated: \$10 billion
Criminals begin using offline wallets to move funds with no transaction monitoring.}
\vspace{0.2cm}
\textbf{Question: What AML risk does offline CBDC functionality present?}
\vspace{0.1cm}
\begin{itemize}
    \item \option{A)} No risk; all CBDC transactions are traceable
    \item \option{B)} Offline transactions cannot be monitored in real time, creating an anonymous payment channel similar to cash but with easier bulk transport
    \item \option{C)} Only the \$200 transaction limit is a concern
    \item \option{D)} Offline CBDC is impossible to use for crime
    \item \option{E)} Only the offline balance limit matters
\end{itemize}
\vspace{0.2cm}
\trap{CBDC is fully traceable; offline changes nothing.}
\correct{Answer: B. Offline CBDC transactions cannot be monitored in real time. While online CBDC provides transparency, offline functionality creates an anonymous payment channel that could be exploited. Criminals could use multiple offline wallets to move funds without monitoring. The \$1,000 balance limit is a control, but large numbers of wallets could be used.}
\redflag{Offline CBDC = monitoring gap. Balance and transaction limits are critical controls.}
\end{frame}

\begin{frame}{DOMAIN B - QUESTION 22: Crypto Airdrop ML Risks}
\scenario{A new cryptocurrency project conducts an "airdrop" distributing 1 million tokens to 100,000 wallet addresses. The airdrop has the following characteristics:
- No purchase required; tokens are "free"
- Recipients only need to provide a wallet address (no KYC)
- Tokens are immediately tradeable on a DEX after distribution
- Token price reaches \$5 within 2 weeks (total value \$5 million)
- Project team retains 20\% of tokens (200,000 tokens)
- Anonymous team claims to be "decentralized"
- After price peaks, team sells all 200,000 tokens for 2,000 ETH (\$4 million)
- Team then launders ETH through mixers and stablecoin conversions
- 50\% of airdrop recipients are shell wallets controlled by team
The project marketed itself as a "community-driven DeFi protocol."}
\vspace{0.2cm}
\textbf{Question: What money laundering technique does this crypto airdrop scheme represent?}
\vspace{0.1cm}
\begin{itemize}
    \item \option{A)} Legitimate token distribution; no laundering
    \item \option{B)} Pump-and-dump with self-dealing; illicit funds converted to "clean" ETH through artificial token value creation
    \item \option{C)} Traditional structuring
    \item \option{D)} Trade-based money laundering
    \item \option{E)} Cash smuggling
\end{itemize}
\vspace{0.2cm}
\trap{Airdrops are legitimate marketing; nothing suspicious.}
\correct{Answer: B. This is a pump-and-dump scheme used for money laundering. The team: (1) creates artificial token value through hype and wash trading, (2) uses shell wallets to receive airdrop, (3) sells tokens at inflated price, (4) converts to ETH/stablecoins, (5) launders through mixers. Initial illicit funds become "clean" crypto through the scheme.}
\redflag{Crypto airdrops with anonymous teams + wash trading + mixers = money laundering scheme.}
\end{frame}

\begin{frame}{DOMAIN B - QUESTION 23: Staking and Yield Farming ML}
\scenario{A criminal organization uses DeFi staking and yield farming protocols to launder funds:
- Initial funds: 1,000 ETH (\$1.8 million) from phishing scams
- Deposits ETH into staking pool (no KYC required)
- Earns 5\% APY in reward tokens for 6 months
- Reward tokens sold for USDC on DEX
- Withdraws original ETH + rewards (approximately 1,025 ETH)
- Transfers ETH through 10 different DeFi protocols
- Converts to USDC
- Cashes out through licensed VASP
- Total time in DeFi: 8 months
- Claims funds are "legitimate DeFi investment returns" with "proof of stake" documentation
The organization has created fake transaction histories showing clean funds entering the DeFi protocols.}
\vspace{0.2cm}
\textbf{Question: How does DeFi staking facilitate money laundering in this scenario?}
\vspace{0.1cm}
\begin{itemize}
    \item \option{A)} Staking is fully transparent and cannot facilitate laundering
    \item \option{B)} Staking creates legitimate-looking "investment returns" to explain fund origin, provides time-based layering, and generates plausible documentation to support false source of funds claims
    \item \option{C)} Only liquidity pools, not staking, facilitate laundering
    \item \option{D)} Staking requires KYC on all protocols
    \item \option{E)} No laundering possible; staking locks funds
\end{itemize}
\vspace{0.2cm}
\trap{Staking is legitimate investment; no ML risk.}
\correct{Answer: B. Staking facilitates laundering by: (1) providing legitimate-looking "investment returns" as explanation for funds, (2) creating time-based layering (8 months in protocols), (3) generating plausible documentation (staking rewards, APY, transaction history), (4) commingling illicit funds with protocol rewards to obscure source.}
\redflag{DeFi staking + fake documentation + no KYC = major ML vulnerability.}
\end{frame}

\begin{frame}{DOMAIN B - QUESTION 24: MEV (Maximal Extractable Value) Laundering}
\scenario{A sophisticated criminal organization uses Maximal Extractable Value (MEV) bots to launder funds on Ethereum. The organization:
- Controls 15\% of all Ethereum staking validators
- Intercepts pending transactions in the mempool
- Reorders transactions to extract value (sandwich attacks, front-running)
- Launders funds through MEV extraction rather than direct transfers
- Illicit funds (100 ETH) are used to execute MEV strategies
- MEV extracted: 150 ETH over 6 months (50 ETH "profit")
- Funds exit as "legitimate MEV extraction" rather than original illicit source
- MEV extraction profits paid directly to clean wallets
- Original 100 ETH remains in protocol but is now layered
Tax authorities cannot determine which ETH is original illicit funds vs MEV profits.}
\vspace{0.2cm}
\textbf{Question: What AML challenge does MEV extraction present for investigators?}
\vspace{0.1cm}
\begin{itemize}
    \item \option{A)} MEV is fully traceable; no challenge
    \item \option{B)} MEV extraction commingles illicit funds with legitimate protocol profits; investigators cannot distinguish original source from extracted value; profits appear as "legitimate trading"
    \item \option{C)} Only sandwich attacks are problematic
    \item \option{D)} MEV only works on Bitcoin; Ethereum not affected
    \item \option{E)} MEV requires KYC, eliminating ML risk
\end{itemize}
\vspace{0.2cm}
\trap{MEV is complex but traceable; no ML issue.}
\correct{Answer: B. MEV extraction commingles funds. The 150 ETH extracted contains both original illicit funds and "profits." Investigators cannot determine which is which. The MEV extraction process creates plausible deniability and legitimizes the funds as "trading profits" rather than original illicit source.}
\redflag{MEV extraction = advanced layering technique that commingles illicit and legitimate funds.}
\end{frame}

% ============================================================
% DOMAIN C - QUESTIONS 25-36
% ============================================================

\begin{frame}{DOMAIN C - QUESTION 25: MiCA Regulation Scope}
\scenario{A crypto company offers the following services in the EU:
- Cryptocurrency exchange (BTC/ETH/USDT)
- Custodial wallet services
- Crypto lending and borrowing
- Staking as a service
- Initial exchange offerings (IEOs)
The company is currently registered only in one EU member state and has been operating under that license for 3 years. The company wants to expand services to all 27 EU member states without registering separately in each country.}
\vspace{0.2cm}
\textbf{Question: Under MiCA (Markets in Cryptoassets Regulation), what is required for this company to operate across the EU?}
\vspace{0.1cm}
\begin{itemize}
    \item \option{A)} Register separately in each of the 27 EU member states
    \item \option{B)} Obtain a single MiCA license in one EU member state with passporting rights to all EU member states
    \item \option{C)} No license required; crypto is unregulated in the EU
    \item \option{D)} Only register with ESMA directly
    \item \option{E)} Only need license for exchange; other services are exempt
\end{itemize}
\vspace{0.2cm}
\trap{Each EU country has separate crypto regulations.}
\correct{Answer: B. MiCA provides a single EU-wide licensing framework. Authorized in one member state = passporting rights to all EU member states. This harmonizes regulation across the EU and replaces fragmented national approaches.}
\redflag{MiCA = single EU license with passporting rights. Authorized in one country = operate in all 27.}
\end{frame}

\begin{frame}{DOMAIN C - QUESTION 26: MiCA Stablecoin Rules}
\scenario{A company issues a stablecoin (ART - Asset-Referenced Token) pegged to a basket of 5 fiat currencies and commodities. The stablecoin has the following characteristics:
- Market capitalization: €5 billion
- Daily trading volume: €500 million
- Users: 2 million across EU and non-EU countries
- Reserves: held in multiple banks across 8 jurisdictions
- Reserve composition: 60\% EUR, 20\% USD, 10\% gold, 5\% Bitcoin, 5\% corporate bonds
- Issuer: established in a non-EU country
- No EU license currently
The company wants to continue serving EU customers.}
\vspace{0.2cm}
\textbf{Question: Under MiCA, what are the requirements for this stablecoin issuer to operate in the EU?}
\vspace{0.1cm}
\begin{itemize}
    \item \option{A)} No requirements; stablecoins are exempt from MiCA
    \item \option{B)} Must obtain license from an EU member state, maintain adequate liquid reserves, and comply with authorization requirements for ARTs
    \item \option{C)} Only need to register; no reserve requirements
    \item \option{D)} Stablecoins over €1 billion market cap are prohibited in EU
    \item \option{E)} Only commodity-backed stablecoins are regulated
\end{itemize}
\vspace{0.2cm}
\trap{Stablecoins are unregulated; no license needed.}
\correct{Answer: B. Under MiCA, ART issuers must be authorized in an EU member state, maintain sufficient liquid reserves (1:1 backing), have qualified directors with good repute, and implement effective AML programs. The €5 billion market cap triggers additional scrutiny.}
\redflag{Stablecoins (ARTs) require EU authorization. Reserves must be liquid and adequately backed.}
\end{frame}

\begin{frame}{DOMAIN C - QUESTION 27: EU 5AMLD Anonymous Prepaid Card Threshold}
\scenario{A financial institution offers prepaid cards with the following features:
- Card 1: Maximum value €200, can be purchased with cash, no ID required
- Card 2: Maximum value €500, requires email address only, no ID verification
- Card 3: Maximum value €1,000, requires basic ID (name and address only)
- Card 4: Maximum value €250, can be purchased online with cryptocurrency, no ID
- All cards can be reloaded up to maximum value
- Cards can be used internationally
- Estimated annual issuance: 500,000 cards
- Law enforcement reports that 30\% of seized terrorist financing materials include prepaid cards
The institution is located in an EU member state.}
\vspace{0.2cm}
\textbf{Question: Under 5AMLD, which prepaid cards are non-compliant with the anonymous card threshold?}
\vspace{0.1cm}
\begin{itemize}
    \item \option{A)} All cards are compliant because AML applies only to bank accounts
    \item \option{B)} Cards 1, 2, 3, and 4 are all non-compliant because the threshold is €150
    \item \option{C)} Only Card 3 is non-compliant because value exceeds €150
    \item \option{D)} Only Card 4 is non-compliant because purchased with crypto
    \item \option{E)} Only reloadable cards are non-compliant
\end{itemize}
\vspace{0.2cm}
\trap{4AMLD threshold was €250, so €200 is fine.}
\correct{Answer: B. 5AMLD lowered the anonymous prepaid card threshold from €250 to €150. Cards with maximum value exceeding €150 must be linked to an identified customer account. Cards 1 (€200), 2 (€500), 3 (€1,000), and 4 (€250) all exceed €150 and require customer identification.}
\redflag{5AMLD anonymous prepaid card limit = €150. Any card exceeding €150 requires KYC.}
\end{frame}

\begin{frame}{DOMAIN C - QUESTION 28: US AML Act 2020 Crypto Coverage}
\scenario{A cryptocurrency exchange operating in the United States offers the following services:
- Spot trading (BTC/USD, ETH/USD, etc.)
- Margin trading up to 5x leverage
- Crypto derivatives (futures and options)
- Staking services
- Crypto lending
- Non-custodial wallet (separate entity, same owner)
The exchange has:
- Not registered as an MSB with FinCEN
- No AML program
- No SAR filings in 4 years of operation
- Annual trading volume: \$10 billion
- Customer base: 500,000 US persons
The exchange claims it is not a "financial institution" under the Bank Secrecy Act.}
\vspace{0.2cm}
\textbf{Question: Under the AML Act of 2020, is this exchange subject to BSA/AML requirements?}
\vspace{0.1cm}
\begin{itemize}
    \item \option{A)} No, cryptocurrency exchanges are explicitly exempt from the BSA
    \item \option{B)} Yes, cryptocurrency exchanges are considered money services businesses (MSBs) with full BSA/AML requirements including registration, program, and SAR filing
    \item \option{C)} Only spot trading is covered; derivatives and lending are exempt
    \item \option{D)} Only exchanges with fiat on-ramps are covered
    \item \option{E)} Only exchanges trading privacy coins are covered
\end{itemize}
\vspace{0.2cm}
\trap{Crypto is unregulated; no BSA requirements.}
\correct{Answer: B. The AML Act of 2020 "updates existing AML regulations to include cryptocurrency exchanges. They are considered money services businesses and have the same licensing and reporting requirements." This includes registration with FinCEN, AML program, SAR filings, and recordkeeping.}
\redflag{US crypto exchanges = MSBs under AML Act 2020. Full BSA compliance required.}
\end{frame}

\begin{frame}{DOMAIN C - QUESTION 29: OFAC Sanctions and Crypto Mixers}
\scenario{A US-based VASP identifies transactions involving the following:
- Transaction A: 5 ETH from a wallet address that interacted with Tornado Cash (OFAC-sanctioned mixer)
- Transaction B: 2 BTC from a wallet address that used ChipMixer (not specifically sanctioned by OFAC)
- Transaction C: 10,000 USDC from a wallet address linked to a North Korean hacking group (OFAC-sanctioned entity)
- Transaction D: 1 BTC from a wallet that used Wasabi Wallet's CoinJoin (not sanctioned)
The VASP has an OFAC compliance program but has not blocked any crypto transactions to date.}
\vspace{0.2cm}
\textbf{Question: Which transactions require blocking under OFAC sanctions?}
\vspace{0.1cm}
\begin{itemize}
    \item \option{A)} None; OFAC does not apply to crypto
    \item \option{B)} Only Transaction C (North Korean hacking group) requires blocking
    \item \option{C)} Transactions A, B, C, and D all require blocking
    \item \option{D)} Only Transaction A (Tornado Cash) requires blocking
    \item \option{E)} Transactions A and C require blocking; B and D require EDD
\end{itemize}
\vspace{0.2cm}
\trap{Only sanctioned entities, not mixers, require blocking.}
\correct{Answer: E. OFAC sanctions apply to specific designated persons and entities. Tornado Cash is specifically sanctioned by OFAC (Transaction A). The North Korean hacking group is a sanctioned entity (Transaction C). ChipMixer and Wasabi Wallet/CoinJoin are not specifically sanctioned (as of this writing), but funds from mixers require EDD.}
\redflag{OFAC sanctions apply to crypto addresses designated on SDN List. Mixers may or may not be sanctioned.}
\end{frame}

\begin{frame}{DOMAIN C - QUESTION 30: FATF Recommendation 15 New Technologies}
\scenario{A bank plans to launch a new cryptocurrency custody service for customers. The service will allow customers to:
- Buy, sell, and hold Bitcoin and Ethereum through the bank's mobile app
- Transfer crypto to external wallets
- Use crypto as collateral for loans
The bank's product development timeline:
- Month 1-2: Requirements gathering
- Month 3-4: Technical development
- Month 5: Internal testing
- Month 6: Soft launch with 100 customers
- Month 7: Full launch
The AML risk assessment for the new crypto service is scheduled for Month 8 (after full launch). No risk assessment was conducted before development began. The bank's CCO was informed of the new service in Month 6.}
\vspace{0.2cm}
\textbf{Question: Under FATF Recommendation 15 (New Technologies), is the bank's approach compliant?}
\vspace{0.1cm}
\begin{itemize}
    \item \option{A)} Yes, risk assessment after launch is acceptable
    \item \option{B)} No, risk assessment must be completed BEFORE launch of new products or technologies
    \item \option{C)} Only need risk assessment if crypto volume exceeds \$1 million
    \item \option{D)} CCO involvement at Month 6 is sufficient
    \item \option{E)} New products are exempt from risk assessment
\end{itemize}
\vspace{0.2cm}
\trap{Risk assessment can be done at any time before, during, or after launch.}
\correct{Answer: B. FATF Recommendation 15 requires financial institutions to "identify and assess the money laundering or terrorist financing risks that may arise in relation to the development of new products and new business practices, including new delivery mechanisms, and the use of new or developing technologies for both new and pre-existing products." This assessment must be done PRIOR to launch.}
\redflag{FATF Rec. 15: Risk assessment for new products/technologies MUST be before launch.}
\end{frame}

\begin{frame}{DOMAIN C - QUESTION 31: EU AMLA Direct Supervision Powers}
\scenario{The EU AML Authority (AMLA) begins operations in Frankfurt. AMLA identifies a cross-border crypto exchange that operates in 15 EU member states with the following characteristics:
- Annual trading volume: €50 billion
- Customer base: 10 million EU residents
- Previously supervised by 15 different national authorities
- National regulators issued conflicting guidance
- Exchange had 5 different AML program interpretations across countries
- Exchange identified as high residual risk by AMLA risk assessment
- Exchange has had 3 major AML compliance failures in past 2 years
- National regulators failed to take consistent enforcement action}
\vspace{0.2cm}
\textbf{Question: Under the EU AML Authority (AMLA) framework, what supervisory powers does AMLA have over this crypto exchange?}
\vspace{0.1cm}
\begin{itemize}
    \item \option{A)} No powers; AMLA only coordinates, does not supervise
    \item \option{B)} Direct supervision authority for high-risk cross-border entities including on-site inspections, corrective measures, and penalties
    \item \option{C)} Only can issue guidance; no enforcement powers
    \item \option{D)} Only can supervise banks, not crypto exchanges
    \item \option{E)} Only can supervise national regulators, not directly supervise entities
\end{itemize}
\vspace{0.2cm}
\trap{AMLA only coordinates; direct supervision remains with national authorities.}
\correct{Answer: B. AMLA has direct supervision authority over selected high-risk cross-border financial institutions, including crypto exchanges. Powers include: on-site inspections, joint supervisory teams, corrective measures, and imposing penalties for serious, repeated, or systematic breaches. AMLA can directly supervise up to 40 high-risk entities.}
\redflag{AMLA = direct supervision of high-risk cross-border entities, including VASPs.}
\end{frame}

\begin{frame}{DOMAIN C - QUESTION 32: FATF Grey List Impact on Crypto}
\scenario{A jurisdiction is placed on the FATF Grey List (jurisdictions under increased monitoring) due to strategic AML/CFT deficiencies, specifically related to cryptoasset regulation. The jurisdiction has:
- No VASP licensing regime
- No Travel Rule implementation
- No crypto transaction monitoring
- No SAR filing requirements for crypto
- 50+ VASPs operating from the jurisdiction
- Crypto trading volume: \$100 billion annually
- Major global crypto exchanges have subsidiaries in this jurisdiction
A correspondent bank in Country X (not on Grey List) has relationships with banks in the Grey List jurisdiction. A VASP incorporated in the Grey List jurisdiction has a correspondent account with a Country X bank.}
\vspace{0.2cm}
\textbf{Question: What actions should financial institutions take regarding crypto-related entities in Grey List jurisdictions?}
\vspace{0.1cm}
\begin{itemize}
    \item \option{A)} No action; Grey List is advisory only
    \item \option{B)} Apply enhanced due diligence to transactions and relationships involving Grey List jurisdictions, including crypto entities
    \item \option{C)} Terminate all relationships with Grey List jurisdictions immediately
    \item \option{D)} Only apply EDD if jurisdiction is also on Black List
    \item \option{E)} Only banks, not crypto entities, need EDD
\end{itemize}
\vspace{0.2cm}
\trap{Grey List = no action needed; only Black List matters.}
\correct{Answer: B. FATF calls on all members to apply enhanced due diligence to Grey List jurisdictions. This includes crypto entities. Countermeasures may be required for persistent deficiencies. Termination is not required but EDD is mandatory.}
\redflag{Grey List = EDD required. Black List = countermeasures potentially including termination.}
\end{frame}

\begin{frame}{DOMAIN C - QUESTION 33: FinCEN 314(b) Crypto Information Sharing}
\scenario{Three VASPs (VASP A, VASP B, VASP C) are all registered MSBs with FinCEN and have opted into 314(b) information sharing. They identify a common customer who has:
- Deposited 500 BTC to VASP A over 30 days from mixed wallets
- Withdrawn 480 BTC from VASP A to a private wallet
- Deposited 450 BTC to VASP B from a different private wallet
- Withdrawn 430 BTC from VASP B to another private wallet
- Deposited 400 BTC to VASP C from another private wallet
- Total volume: approximately \$15 million
- Customer's declared business: small retail store with \$200k annual revenue
- Each VASP individually sees suspicious activity but limited data
VASP A wants to share information with VASP B and VASP C to determine if this is a coordinated money laundering scheme.}
\vspace{0.2cm}
\textbf{Question: Under FinCEN 314(b), can these VASPs share customer information, and what is required?}
\vspace{0.1cm}
\begin{itemize}
    \item \option{A)} No, 314(b) does not apply to VASPs; only banks
    \item \option{B)} Yes, VASPs are eligible institutions; they must have filed opt-in notices and may share information regarding suspected ML/TF
    \item \option{C)} Yes, but only if the customer consents
    \item \option{D)} No, 314(b) sharing is prohibited for crypto
    \item \option{E)} Yes, but only for fiat transactions, not crypto
\end{itemize}
\vspace{0.2cm}
\trap{314(b) is only for banks, not crypto.}
\correct{Answer: B. VASPs that are registered MSBs are eligible for 314(b). Requirements: (1) file opt-in notice with FinCEN, (2) maintain confidentiality, (3) sharing limited to suspected ML/TF, (4) safe harbor applies. The three VASPs can share information about the common customer's activity across platforms.}
\redflag{314(b) applies to MSBs including VASPs. Opt-in required before sharing.}
\end{frame}

\begin{frame}{DOMAIN C - QUESTION 34: UNSCR 1373 vs 1267 Crypto Asset Freezes}
\scenario{A bank identifies a customer who sends \$50,000 in Bitcoin monthly to a wallet address. The bank's investigation reveals:
- The wallet address is not on any UN Sanctions List (1267/1988)
- The customer's country has domestically designated the recipient as a terrorist financier under UNSCR 1373 implementing legislation
- No court order has been issued
- The recipient is a known associate of a designated terrorist organization
- The customer claims the funds are "charitable donations"
- The bank's home country has implemented UNSCR 1373 into national law}
\vspace{0.2cm}
\textbf{Question: Under UNSCR 1373 vs. UNSCR 1267/1988, must the bank freeze the crypto assets?}
\vspace{0.1cm}
\begin{itemize}
    \item \option{A)} No, only UN 1267/1988 designations require freezing
    \item \option{B)} Yes, UNSCR 1373 requires countries to freeze assets of domestically designated terrorist financiers; bank must comply with national law implementing 1373
    \item \option{C)} Only if amount exceeds \$100,000
    \item \option{D)} Only fiat assets, not crypto, must be frozen
    \item \option{E)} Only if a court order is obtained
\end{itemize}
\vspace{0.2cm}
\trap{No UN listing = no freezing obligation.}
\correct{Answer: B. UNSCR 1373 requires member states to "freeze without delay funds and other financial assets or economic resources of persons who commit, or attempt to commit, terrorist acts." Countries may create their own designation lists under 1373. The bank must comply with national law implementing 1373, even without UN listing.}
\redflag{UNSCR 1373 = countries can create own TF designation lists. UNSCR 1267/1988 = UN Security Council list. Both require asset freezes.}
\end{frame}

\begin{frame}{DOMAIN C - QUESTION 35: EU 7AMLD Cash Payment Limit and Crypto}
\scenario{Under the EU 7AMLD, a new €10,000 cash payment limit is implemented across all EU member states. A customer in France wants to purchase cryptocurrency through a physical crypto ATM. The customer attempts to:
- Transaction 1: Insert €9,500 in cash to buy Bitcoin
- Transaction 2: Insert €9,500 in cash to buy Ethereum (different ATM, same day)
- Transaction 3: Insert €9,500 in cash to buy USDT (next day)
- Transaction 4: Insert €10,500 in cash to buy Bitcoin (single transaction)
- Transaction 5: €8,000 cash for Bitcoin, plus €3,000 wire transfer for same purchase
The crypto ATM operator has no KYC for transactions under €10,000.}
\vspace{0.2cm}
\textbf{Question: Under 7AMLD, which transactions violate the €10,000 cash payment limit?}
\vspace{0.1cm}
\begin{itemize}
    \item \option{A)} None; crypto ATMs are exempt from cash payment limits
    \item \option{B)} Only Transaction 4 (€10,500) violates the limit
    \item \option{C)} Transactions 1, 2, 3, 4, and 5 all violate the limit
    \item \option{D)} Only Transactions 1, 2, and 3 (structuring) violate the limit
    \item \option{E)} Only Transaction 5 (mixed payment) violates the limit
\end{itemize}
\vspace{0.2cm}
\trap{Each transaction under €10,000 is fine; structuring detection is separate.}
\correct{Answer: C. 7AMLD prohibits cash payments exceeding €10,000 for ANY commercial transaction, including crypto purchases. Transaction 4 exceeds €10,000. Transactions 1,2,3 are structuring (multiple transactions under threshold but same customer/related transactions). Transaction 5 combines cash + wire; the cash portion €8,000 is under threshold but must be considered with other transactions. The operator's "no KYC under €10,000" policy is non-compliant with 7AMLD's structuring detection requirements.}
\redflag{7AMLD €10k cash limit + structuring detection required. Cannot avoid by breaking into smaller transactions.}
\end{frame}

\begin{frame}{DOMAIN C - QUESTION 36: DNFBP Crypto Obligations under FATF}
\scenario{A law firm provides the following services to crypto clients:
- Forms legal entities (corporations, LLCs, trusts) for crypto businesses
- Provides registered agent services
- Assists with VASP license applications
- Provides legal advice on crypto regulatory compliance
- Holds client funds in escrow for crypto acquisitions
- Facilitates introductions to banks for crypto clients
The law firm does not directly handle cryptocurrency; all transactions are in fiat currency. The firm claims it is not subject to AML/CFT obligations because it is a legal professional and not a VASP or financial institution.}
\vspace{0.2cm}
\textbf{Question: Under FATF Recommendation 23 (DNFBPs), does this law firm have AML/CFT obligations related to its crypto client services?}
\vspace{0.1cm}
\begin{itemize}
    \item \option{A)} No, law firms are exempt from all AML obligations
    \item \option{B)} Yes, when acting as formation agent, registered agent, or holding client funds, the firm has CDD and reporting obligations
    \item \option{C)} Only when directly handling cryptocurrency
    \item \option{D)} Only for clients that are sanctioned
    \item \option{E)} Only for clients with over €1 million in crypto
\end{itemize}
\vspace{0.2cm}
\trap{Law firms have attorney-client privilege that exempts them from AML.}
\correct{Answer: B. Under FATF Rec. 23, DNFBPs including lawyers have AML obligations when preparing for or carrying out transactions concerning: (1) buying/selling real estate, (2) managing client funds/securities, (3) managing bank/savings accounts, (4) organizing contributions for company creation, (5) creating/operating legal persons or arrangements. Forming legal entities, registered agent services, and holding client funds are covered activities.}
\redflag{Lawyers = DNFBPs with AML obligations for specified activities, including company formation.}
\end{frame}

% ============================================================
% DOMAIN D - QUESTIONS 37-45
% ============================================================

\begin{frame}{DOMAIN D - QUESTION 37: Blockchain Explorer Data Reliability}
\scenario{An AML investigator uses three different blockchain explorers to trace funds from a ransomware attack:
- Explorer A shows transaction from Wallet X to Wallet Y at 14:32 UTC, amount 50 BTC
- Explorer B shows the same transaction at 14:35 UTC, amount 49.99 BTC (after fee deduction)
- Explorer C shows the transaction but lists different "from" address due to address reuse labeling
- Explorer A also incorrectly labels a known exchange wallet as "suspicious - unlicensed"
- Explorer B fails to show 5\% of transactions due to API rate limiting
- Explorer C shows correct data for all transactions but charges \$500/month
The investigator needs to produce evidence for court.}
\vspace{0.2cm}
\textbf{Question: What is the best practice for using blockchain explorers in financial crime investigations?}
\vspace{0.1cm}
\begin{itemize}
    \item \option{A)} Trust the most popular free explorer
    \item \option{B)} Compare data across multiple explorers, verify reliability, and use paid enterprise tools for court-admissible evidence
    \item \option{C)} Never use explorers; only use law enforcement subpoenas
    \item \option{D)} Accept all explorer data as equally reliable
    \item \option{E)} Only use explorers for preliminary research; all data is unreliable
\end{itemize}
\vspace{0.2cm}
\trap{All blockchain explorers show the same data; any is fine.}
\correct{Answer: B. Blockchain explorers are third-party tools that may present false information. Different explorers may show different data due to indexing differences, labeling errors, or API limitations. Best practice: (1) compare data across multiple explorers, (2) verify reliability, (3) use paid enterprise tools (Chainalysis, Elliptic, TRM Labs) for court-admissible evidence.}
\redflag{Blockchain explorers = potential data discrepancies. Cross-verify and use paid tools for evidence.}
\end{frame}

\begin{frame}{DOMAIN D - QUESTION 38: Clustering Analysis for Crypto Investigation}
\scenario{An investigator analyzes Bitcoin transactions and identifies the following address relationships:
- Address A and Address B are both inputs to Transaction 1
- Address B and Address C are both inputs to Transaction 2
- Address C and Address D are both inputs to Transaction 3
- Address D and Address E are both inputs to Transaction 4
- Address E and Address A are both inputs to Transaction 5
- Address F has never been used as input with any other address
- Address G is always used alone as input
- Transaction 6: Address H and Address I are inputs; Address J and Address K are outputs
- Address H and J appear together in 10 other transactions
The investigator needs to group addresses that likely belong to the same entity.}
\vspace{0.2cm}
\textbf{Question: Which addresses would clustering analysis group together as belonging to the same entity?}
\vspace{0.1cm}
\begin{itemize}
    \item \option{A)} No addresses can be grouped; Bitcoin is anonymous
    \item \option{B)} Addresses A, B, C, D, and E (connected through common input clustering)
    \item \option{C)} Only Address F and G together
    \item \option{D)} Address H, I, J, and K together
    \item \option{E)} All addresses A through K belong to the same entity
\end{itemize}
\vspace{0.2cm}
\trap{Each address is independent; no grouping possible.}
\correct{Answer: B. Clustering uses the common-input heuristic: when multiple addresses are inputs to the same transaction, they are likely controlled by the same entity. A+B in Tx1, B+C in Tx2, C+D in Tx3, D+E in Tx4, E+A in Tx5 creates a cluster containing A,B,C,D,E. Address F,G are isolated; H,I,J,K have patterns suggesting entity relationships but need additional analysis.}
\redflag{Common-input heuristic = multiple addresses as inputs = same entity likely.}
\end{frame}

\begin{frame}{DOMAIN D - QUESTION 39: AI Transaction Monitoring for Crypto}
\scenario{A VASP implements an AI-based transaction monitoring system to replace its legacy rules-based system. The legacy system had:
- 10,000 alerts per month
- 98\% false positive rate
- Missed 30\% of actual suspicious activity (false negative rate 30\%)
- 15 monitoring scenarios, all threshold-based
The new AI system after 6 months of tuning shows:
- 2,000 alerts per month
- 60\% false positive rate (2,000 alerts, 800 true positives, 1,200 false positives)
- 5\% false negative rate
- Detected 3 new typologies that legacy system missed entirely
- Requires retraining every 3 months
The VASP's compliance staff is concerned about the AI system's "black box" nature and inability to explain why specific alerts were generated.}
\vspace{0.2cm}
\textbf{Question: What are the trade-offs between rules-based and AI-based transaction monitoring for crypto?}
\vspace{0.1cm}
\begin{itemize}
    \item \option{A)} AI is always better with no downsides
    \item \option{B)} AI reduces false positives and false negatives but requires explainability and ongoing governance; black-box models may not satisfy regulatory expectations
    \item \option{C)} Rules-based systems are always superior for crypto
    \item \option{D)} AI cannot be used for crypto monitoring
    \item \option{E)} Only false positives matter; false negatives are irrelevant
\end{itemize}
\vspace{0.2cm}
\trap{AI is perfect; no issues with black-box models.}
\correct{Answer: B. AI reduces false positives (98\%→60\%) and false negatives (30\%→5\%) significantly. However, AI presents challenges: (1) explainability - regulators expect to understand why alerts are generated, (2) black-box models may not satisfy regulatory requirements, (3) requires ongoing testing and retraining, (4) potential bias in training data.}
\redflag{AI monitoring = better detection but requires explainability and governance. Black-box models risk regulatory rejection.}
\end{frame}

\begin{frame}{DOMAIN D - QUESTION 40: Crypto ATMs and Geolocation Red Flags}
\scenario{A bank monitors a customer who regularly uses crypto ATMs. Over a 3-month period:
- Customer withdraws \$500 cash from Bank Branch A (9:00 AM)
- Within 30 minutes, \$500 in Bitcoin is purchased from Crypto ATM located 0.5 miles from Branch A
- Pattern repeats 45 times over 3 months
- Customer's home address is 50 miles from all ATM locations
- Customer's work address is in different city
- Customer cannot explain why they travel 50 miles to use specific ATMs
- Crypto ATMs used are from 3 different operators, all with \$900 daily limits without ID
- Customer's declared income: \$60,000/year
- Total crypto purchased: \$22,500 over 3 months
- Customer has no other crypto activity or trading accounts}
\vspace{0.2cm}
\textbf{Question: What geolocation-related red flags are present in this scenario?}
\vspace{0.1cm}
\begin{itemize}
    \item \option{A)} No red flags; crypto ATMs are legal
    \item \option{B)} Geographic distance from home/work, timing correlation with bank withdrawals, use of multiple ATMs with low ID thresholds, and income inconsistency
    \item \option{C)} Only the \$500 withdrawal amount is a red flag
    \item \option{D)} Only the use of multiple ATM operators is a red flag
    \item \option{E)} Only the income inconsistency matters
\end{itemize}
\vspace{0.2cm}
\trap{Crypto ATMs are legal; travel distance is not suspicious.}
\correct{Answer: B. Geolocation red flags include: (1) transactions from locations inconsistent with customer profile (50 miles from home), (2) timing correlation between bank withdrawal and crypto purchase (30 minutes), (3) use of ATMs with low ID thresholds (\$900), (4) income inconsistency (\$60k income vs \$22,500 crypto purchases in 3 months = 150\% of annual income).}
\redflag{Geolocation + timing + income inconsistency = multiple red flags requiring investigation.}
\end{frame}

\begin{frame}{DOMAIN D - QUESTION 41: Device Fingerprinting for Crypto Fraud}
\scenario{A VASP's fraud detection system uses device fingerprinting and identifies the following patterns:
- Device 1 (fingerprint ABC): Used to access 50 different customer accounts over 30 days
- Device 2 (fingerprint DEF): Used to access 1 customer account only, but from 200 different IP addresses across 15 countries in 24 hours
- Device 3 (fingerprint GHI): New device, no history, accessed account with 3 failed login attempts then successful on 4th
- Device 4 (fingerprint JKL): Known fingerprint associated with 3 previous fraud cases
- Device 5 (fingerprint MNO): Customer's registered device for 2 years, normal behavior pattern
The VASP has implemented device fingerprinting but has not integrated it with transaction monitoring.}
\vspace{0.2cm}
\textbf{Question: Which devices present the highest fraud or money laundering risk based on device fingerprinting?}
\vspace{0.1cm}
\begin{itemize}
    \item \option{A)} Device 5 only; all others are normal
    \item \option{B)} Devices 1, 2, 3, and 4 present elevated risk; Device 5 is normal
    \item \option{C)} Only Device 2 presents risk due to IP diversity
    \item \option{D)} Only Device 4 presents risk due to fraud history
    \item \option{E)} Device fingerprinting cannot detect risk
\end{itemize}
\vspace{0.2cm}
\trap{Only known fraud devices are risky.}
\correct{Answer: B. Device 1: 50 accounts on same device = potential mule network or account takeover. Device 2: 200 IPs across 15 countries = VPN/proxy usage or automated access. Device 3: new device with failed logins = potential account takeover attempt. Device 4: known fraud device = high risk. Device 5: normal pattern = low risk. Device fingerprinting should be integrated with transaction monitoring for full detection.}
\redflag{Device fingerprinting red flags: multiple accounts on one device, impossible travel (many IPs), failed logins, known fraud devices.}
\end{frame}

\begin{frame}{DOMAIN D - QUESTION 42: Perpetual KYC for Crypto Customers}
\scenario{A bank implements perpetual KYC (pKYC) for its crypto customer segment. Traditional KYC was conducted annually. Under pKYC, the system triggers reviews based on:
- Transaction volume increase: >200\% of 90-day rolling average
- New jurisdiction exposure: first transaction to/from FATF Grey List country
- Wallet type change: new wallet addresses from mixers or privacy coins
- Velocity changes: >10 transactions per day (previous average 2 per day)
- Counterparty risk: receiving funds from VASP on regulatory watchlist
- Static data changes: address, phone, email changes
- Adverse media: new negative news alerts
The system triggers 1,500 reviews in first month, of which 300 require EDD and 50 result in SAR filings. Under traditional annual KYC, these 1,500 triggers would have been missed for up to 11 months.}
\vspace{0.2cm}
\textbf{Question: What are the advantages of perpetual KYC over traditional periodic KYC for crypto customers?}
\vspace{0.1cm}
\begin{itemize}
    \item \option{A)} No advantages; periodic KYC is sufficient
    \item \option{B)} Real-time detection of risk changes, immediate trigger-based reviews, reduced window of undetected suspicious activity, and efficient resource allocation to highest-risk triggers
    \item \option{C)} Only reduces paperwork
    \item \option{D)} Only works for traditional banking, not crypto
    \item \option{E)} pKYC eliminates need for any customer due diligence
\end{itemize}
\vspace{0.2cm}
\trap{Annual KYC is sufficient for all customers regardless of activity.}
\correct{Answer: B. Traditional periodic KYC (annual) leaves up to 11 months of undetected risk changes. Perpetual KYC: (1) real-time detection of triggers, (2) immediate reviews when risk changes, (3) eliminates window between reviews, (4) resources focused on active risks rather than all customers on a schedule. Critical for volatile crypto customer behavior.}
\redflag{Perpetual KYC = real-time risk detection. Critical for crypto where risk changes rapidly.}
\end{frame}

\begin{frame}{DOMAIN D - QUESTION 43: Open Source Intelligence (OSINT) for Crypto Investigations}
\scenario{An AML investigator is researching a crypto wallet address linked to suspected ransomware activity. Using OSINT techniques, the investigator:
- Search 1: Wallet address in blockchain explorer → shows transactions
- Search 2: Wallet address in Google → finds forum post where user posted address for "help with transaction"
- Search 3: Username from forum post in Twitter → finds same username discussing crypto trading
- Search 4: Twitter profile reveals real name and location
- Search 5: Real name in LinkedIn → finds employer in IT security
- Search 6: Employer in corporate registry → finds company director name
- Search 7: Director name in sanctions lists → no match
- Search 8: Director name in PEP database → no match
- Search 9: Director name in adverse media → previous fraud conviction
The investigator builds a profile linking the wallet address to a real person with a fraud conviction.}
\vspace{0.2cm}
\textbf{Question: What OSINT techniques were used in this investigation, and what are their limitations?}
\vspace{0.1cm}
\begin{itemize}
    \item \option{A)} OSINT is not reliable for crypto investigations
    \item \option{B)} Blockchain explorer, search engine, social media, corporate registry, and sanctions/PEP/adverse media checks; limitations include data accuracy, source reliability, and potential false associations
    \item \option{C)} Only blockchain explorers are useful
    \item \option{D)} Only social media is useful
    \item \option{E)} OSINT always produces 100\% accurate results
\end{itemize}
\vspace{0.2cm}
\trap{OSINT is always accurate and reliable.}
\correct{Answer: B. OSINT techniques used: blockchain explorers (transaction data), search engines (forum posts), social media (Twitter, LinkedIn), corporate registries (company director), and sanctions/PEP/adverse media databases. Limitations: (1) source reliability varies, (2) information may be outdated, (3) false associations possible, (4) requires cross-verification across multiple sources.}
\redflag{OSINT = valuable but requires source verification. Cross-reference multiple sources.}
\end{frame}

\begin{frame}{DOMAIN D - QUESTION 44: Entity Resolution for Crypto Corporate Structures}
\scenario{A bank uses entity resolution technology to analyze a crypto trading company customer. The system analyzes:
- Corporate registry data: Company A registered to Address 1, Director John Smith
- Utility bill: Address 1 listed for electricity, account in name "JS Holdings"
- Transaction data: Company A sends funds to Company B
- Corporate registry: Company B registered to Address 1 (same), Director Jane Smith (John's spouse)
- Social media: John Smith's LinkedIn lists himself as "Director, JS Holdings"
- Bank account: Company A and Company B both list same phone number
- Crypto exchange account: Company A's exchange account linked to wallet X
- Wallet X: Also used to fund Company B's exchange account
- Adverse media: John Smith mentioned in crypto fraud investigation
The entity resolution system links these data points together despite different legal names.}
\vspace{0.2cm}
\textbf{Question: What does entity resolution reveal about the relationship between Company A and Company B?}
\vspace{0.1cm}
\begin{itemize}
    \item \option{A)} Company A and Company B are completely unrelated
    \item \option{B)} Entity resolution reveals common beneficial ownership (John Smith), shared address, shared phone number, and linked wallets; likely related entities requiring consolidated risk assessment
    \item \option{C)} Only the shared address is relevant
    \item \option{D)} Only the shared phone number matters
    \item \option{E)} Entity resolution cannot link companies
\end{itemize}
\vspace{0.2cm}
\trap{Separate legal entities are independent; no linkage needed.}
\correct{Answer: B. Entity resolution links: common address (Address 1), common phone number, common individual (John Smith connected to both through directorship and JS Holdings), linked wallets (same wallet X used for both companies' exchange accounts), and adverse media (John Smith fraud investigation). The entities should be treated as related for risk assessment, with aggregated exposure considered.}
\redflag{Entity resolution = links separate legal entities with common control. Critical for crypto corporate structures.}
\end{frame}

\begin{frame}{DOMAIN D - QUESTION 45: RPA for Crypto SAR Filing}
\scenario{A VASP processes 50,000 suspicious activity alerts per month, resulting in 500 SAR filings. The manual SAR filing process requires:
- Investigator to copy data from 6 different systems
- Manually populate 45 form fields
- Write narrative summarizing investigation (30-60 minutes per SAR)
- Supervisor review and approval (15 minutes per SAR)
- Compliance officer final review (10 minutes per SAR)
- Total time per SAR: 55-85 minutes
- Monthly SAR filing time: 458-708 hours
The VASP implements Robotic Process Automation (RPA) to:
- Automatically extract data from 6 systems
- Auto-populate 40 of 45 form fields
- Generate draft narrative using templates based on alert type
- Flag incomplete data for manual completion
- Route to supervisor and compliance officer
Results after implementation:
- Time per SAR: 15-25 minutes (70\% reduction)
- Data entry errors: 95\% reduction
- SAR filing backlog: eliminated
- Staff redeployed to complex investigations}
\vspace{0.2cm}
\textbf{Question: What are the benefits and risks of using RPA for SAR filing in crypto AML compliance?}
\vspace{0.1cm}
\begin{itemize}
    \item \option{A)} No benefits; RPA cannot be used for SAR filing
    \item \option{B)} Benefits: efficiency, reduced errors, faster filing, resource optimization; Risks: automation of flawed processes, over-reliance without human review, template narratives may miss unique red flags
    \item \option{C)} Only risks; no benefits
    \item \option{D)} RPA eliminates need for human investigators entirely
    \item \option{E)} Only benefits; no risks
\end{itemize}
\vspace{0.2cm}
\trap{RPA is perfect; no risks.}
\correct{Answer: B. Benefits: (1) 70\% time reduction, (2) 95\% error reduction, (3) backlog elimination, (4) staff focused on complex cases. Risks: (1) automating flawed processes without fixing underlying issues, (2) over-reliance on automation, (3) template narratives may miss unique circumstances, (4) human review still required for quality assurance.}
\redflag{RPA = efficiency gains but requires process reengineering first. Human review still critical.}
\end{frame}

% ============================================================
% ANSWER KEY SUMMARY SLIDE
% ============================================================

\begin{frame}{ANSWER KEY SUMMARY}
\begin{multicols}{2}
\begin{itemize}
    \item QA-1: \correct{E}
    \item QA-2: \correct{C}
    \item QA-3: \correct{B}
    \item QA-4: \correct{B}
    \item QA-5: \correct{B}
    \item QA-6: \correct{B}
    \item QA-7: \correct{C}
    \item QA-8: \correct{B}
    \item QA-9: \correct{B}
    \item QA-10: \correct{B}
    \item QA-11: \correct{B}
    \item QA-12: \correct{B}
    \item QB-13: \correct{B}
    \item QB-14: \correct{C}
    \item QB-15: \correct{B}
    \item QB-16: \correct{B}
    \item QB-17: \correct{B}
    \item QB-18: \correct{B}
    \item QB-19: \correct{B}
    \item QB-20: \correct{B}
    \item QB-21: \correct{B}
    \item QB-22: \correct{B}
    \item QB-23: \correct{B}
    \item QB-24: \correct{B}
    \item QC-25: \correct{B}
    \item QC-26: \correct{B}
    \item QC-27: \correct{B}
    \item QC-28: \correct{B}
    \item QC-29: \correct{E}
    \item QC-30: \correct{B}
    \item QC-31: \correct{B}
    \item QC-32: \correct{B}
    \item QC-33: \correct{B}
    \item QC-34: \correct{B}
    \item QC-35: \correct{C}
    \item QC-36: \correct{B}
    \item QD-37: \correct{B}
    \item QD-38: \correct{B}
    \item QD-39: \correct{B}
    \item QD-40: \correct{B}
    \item QD-41: \correct{B}
    \item QD-42: \correct{B}
    \item QD-43: \correct{B}
    \item QD-44: \correct{B}
    \item QD-45: \correct{B}
\end{itemize}
\end{multicols}
\vspace{0.3cm}
\centering
\greennote{45 Hard Questions | All 4 Domains | Bitcoin, Crypto, VASPs, DeFi, MiCA, Travel Rule}
\end{frame}

\begin{frame}{Final Summary - Key Crypto AML Concepts Tested}
\begin{itemize}
    \item \textbf{Microstructuring:} Small crypto transactions to avoid detection (Domain A)
    \item \textbf{Mixers/Tumblers:} CoinJoin, Wasabi, Tornado Cash - make tracing nearly impossible
    \item \textbf{Privacy Coins:} Monero, Zcash - higher risk than transparent blockchains
    \item \textbf{Stablecoins:} Fiat-collateralized = ease of fiat conversion = ML risk
    \item \textbf{NFT Wash Trading:} Self-laundering through overpriced NFT sales
    \item \textbf{DeFi VASP Status:} Control/influence test - governance tokens + admin keys
    \item \textbf{Travel Rule:} Originator/beneficiary info required for VASPs same as banks
    \item \textbf{MiCA:} Single EU license, passporting rights, €10k cash limit
    \item \textbf{US AML Act:} Crypto exchanges = MSBs with full BSA requirements
    \item \textbf{FATF Rec. 15:} New product risk assessment BEFORE launch
    \item \textbf{Blockchain Tracing:} On-chain + off-chain analytics to identify real-world persons
    \item \textbf{Clustering:} Common-input heuristic groups addresses by entity
    \item \textbf{Perpetual KYC:} Real-time trigger-based reviews for crypto
    \item \textbf{OSINT:} Cross-source verification critical; limitations exist
    \item \textbf{RPA:} Efficiency gains but requires process reengineering first
\end{itemize}
\vspace{0.3cm}
\centering
\greennote{Good luck on your CAMS exam!}
\end{frame}

\end{document}